\section{Benchmark Setup}
\label{sec:setup}

This section specifies the datasets, the panel of tools under evaluation, the
seven-step harmonization pipeline that places heterogeneous predictors on a
common footing, the per-patient evaluation protocol, and the metrics together
with the fairness safeguards that govern every comparison reported in
Section~\ref{sec:results}. The design goal throughout is a single,
\emph{apples-to-apples} measurement: each tool sees the same inputs, is reduced
to the same peptide-level score under the same aggregation rule, and is scored
against the same continuous ground truth under the same thresholds.

\subsection{Datasets}
\label{sec:datasets}

We evaluate on two cohorts of neoantigen candidates with functional T-cell
readouts measured by interferon-$\gamma$ ELISpot. ELISpot spot-forming counts
(SFC) provide a quantitative, continuous measure of response magnitude rather
than a binary recognition call, and have been used as a quantitative
immunogenicity readout in prior work~\cite{TESLA2020}. The first cohort (DS1)
serves as a development set on which the harmonization pipeline and its scoring
conventions were established and reproduced. DS1 contains $82$ peptides
($325$ sub-peptide rows) and is entirely positive (ELISpot SFC ranging from $16$
upward, with no negative controls); it therefore does not admit a binary
discrimination analysis and is used only to check rank-correlation behavior, where
all evaluated tools were non-significant ($|\rho| \le 0.16$). The second cohort
(DS2) is the primary benchmark and is reported in full below.

DS2 comprises $101$ peptides. Because the tools accept inputs of differing
lengths and allele specificities, each peptide is expanded by exhaustive
sliding-window enumeration over its constituent sub-peptides paired with the
relevant HLA alleles, yielding $33{,}922$ sub-peptide~$\times$~HLA rows
(Section~\ref{sec:harmonization}). The ELISpot SFC value is recorded once per
peptide and propagated to all of its rows. We treat SFC both as a continuous
magnitude target and, after thresholding, as a binary positive/negative label.
Three thresholds are used: SFC $>0$ ($n_{\mathrm{pos}}=90$,
$n_{\mathrm{neg}}=11$), SFC $>10$ ($n_{\mathrm{pos}}=64$), and SFC above the
cohort median of $22.67$ ($n_{\mathrm{pos}}=50$). The strong class imbalance at
the $>0$ threshold (positivity $\approx 89\%$) is intrinsic to this cohort and
directly motivates the use of precision--recall--based metrics
(Section~\ref{sec:metrics}).

\subsection{The tools surveyed}
\label{sec:tools}

The survey targets a panel of $30$ tools spanning the biological cascade from
peptide--HLA binding to functional T-cell immunogenicity: $10$
presentation/binding predictors and $20$ immunogenicity predictors. We organize
this panel into five categories by the quantity each tool natively predicts:
(1)~binding affinity, (2)~presentation / eluted-ligand likelihood,
(3)~immunogenicity via machine-learning or deep models, (4)~immunogenicity via
statistical / sequence-based scores, and (5)~immunogenicity via structural,
foreignness, or composite-pipeline scores. Table~\ref{tab:roster} lists the full
roster together with, for every tool, its output score name, native prediction
task, whether it produces scores for both the mutant (MT) and wild-type (WT)
peptide (a prerequisite for the differential-agretopicity transforms of
Section~\ref{sec:harmonization}), the peptide-length window it scores, whether it
is HLA-aware or HLA-agnostic, its license, and its current integration status.

\paragraph{Integration status (honesty about coverage).} The $30$-tool panel is a
\emph{target}. At the time of writing, $21$ tools are integrated and benchmarked,
producing $22$ harmonized score columns (Status \checkmark in
Table~\ref{tab:roster})---MHCflurry contributes two columns from a single model
(presentation and affinity), whereas BigMHC's eluted-ligand (EL) and
immunogenicity (IM) columns come from two distinct models. A further eight tools
are in progress ($\circ$), and five additional candidates are blocked or have been
excluded ($\times$); the latter are listed only for transparency and are
\emph{not} counted among the $30$. We do not force the integrated and in-progress
counts into an exact arithmetic identity with the $30$ target, because some
in-progress entries (for example, the additional netMHCpan Aff/EL columns) extend
an already-integrated tool rather than adding a new one; the panel should be read
as a $30$-tool target of which $21$ are currently integrated and the remainder are
in progress. Two of the integrated tools are reported with caveats and are never
ranked apples-to-apples against the rest (marked $\ddagger$): HLAthena is included
only as a \emph{presentation proxy}, because peptide presentation is a necessary
but not sufficient condition for a functional T-cell response and its near-random
performance serves as a calibration anchor rather than a competitor; and NetTepi
covers only a subset of the cohort's alleles (no HLA-C), so its column is sparse.
We never present an in-progress or blocked tool as if it were benchmarked. Among
the integrated tools, the eight immunogenicity predictors with full per-peptide
coverage---DeepImmuno, PredIG, IMPROVE, NeoTImmuML$^\star$, pTuneos, PRIME,
ImmuneApp, and deepHLApan---form the primary apples-to-apples panel analyzed in
detail in Section~\ref{sec:results}; their per-tool coverage and missingness are
given in Table~\ref{tab:coverage}.

\begin{table}[t]
\centering
\caption{The $30$-tool target roster surveyed by QuantImmu, grouped into five
categories. Columns: output score name; native prediction task; whether scores
are produced for both mutant and wild-type peptides (MT/WT); native
peptide-length window; HLA-aware vs.\ HLA-agnostic (\emph{agn.}; \emph{partial}
= position-based with allele-specific masks for a subset of alleles); license
(acad.\ = academic-use; n.c.\ = non-commercial); and integration status.
\textbf{Status:} \checkmark\ = integrated and benchmarked; $\circ$ = integration
in progress; $\times$ = blocked or excluded (not part of the $30$). \pendingDTU\
marks tools under a DTU academic license whose benchmark numbers may
not be redistributed to third parties without written consent (Section~\ref{sec:fairness});
$\star$ denotes a self-retrained model (official weights unavailable); $\ddagger$
denotes an integrated tool reported with a caveat and never ranked
apples-to-apples (HLAthena, presentation proxy; NetTepi, low allele coverage).}
\label{tab:roster}
\scriptsize
\setlength{\tabcolsep}{3pt}
\resizebox{\textwidth}{!}{%
\begin{tabular}{llllccll}
\toprule
Tool & Output score & Native task & MT/WT & Len. & HLA & License & Status \\
\midrule
\multicolumn{8}{l}{\textit{(1) Binding affinity}}\\
netMHCpan-4.1 (BA)~\cite{NetMHCpan2020}      & netmhcpan\_ba       & affinity ($-$nM)        & both & 8--14 & aware & DTU\pendingDTU        & \checkmark \\
netMHCpan-4.1 (Aff)~\cite{NetMHCpan2020}     & netMHCpan\_Aff      & affinity ($-$nM)        & both & 8--14 & aware & DTU\pendingDTU        & $\circ$ \\
MHCflurry-2.0 (aff.)~\cite{MHCflurry2020}    & MHCflurry\_affinity & affinity ($-$nM)        & both & 8--15 & aware & Apache-2.0            & \checkmark \\
MHCnuggets~\todo{cite MHCnuggets}            & MHCnuggets ($-$IC$_{50}$) & affinity          & MT   & 8--15 & aware & BSD-3                & \checkmark \\
\midrule
\multicolumn{8}{l}{\textit{(2) Presentation / eluted-ligand}}\\
MHCflurry-2.0 (pres.)~\cite{MHCflurry2020}   & MHCflurry\_pres.    & presentation prob.      & both & 8--15 & aware & Apache-2.0           & \checkmark \\
netMHCpan-4.1 (EL)~\cite{NetMHCpan2020}      & netMHCpan\_EL       & EL \%rank               & both & 8--14 & aware & DTU\pendingDTU       & $\circ$ \\
BigMHC\_EL~\cite{BigMHC2023}                 & BigMHC\_EL          & EL prob.                & both & 8--15 & aware & acad.\ (n.c.)        & \checkmark \\
HLAthena~\cite{HLAthena2020}                 & HLAthena            & presentation prob.      & MT   & 8--11 & aware & academic             & \checkmark$^{\ddagger}$ \\
NetMHCstabpan-1.0~\cite{NetMHCstabpan2016}   & stab.\ \%rank       & stability               & both & 8--14 & aware & DTU\pendingDTU       & $\circ$ \\
TransHLA~\todo{cite TransHLA}                & MT\_TransHLA        & epitope prob.           & MT   & 8--14 & \emph{agn.} & MIT            & \checkmark \\
MAAP~\todo{identify/cite MAAP}               & ---                 & presentation            & ---  & ---   & ---   & ---                  & $\times$ \\
\midrule
\multicolumn{8}{l}{\textit{(3) Immunogenicity --- ML / deep}}\\
DeepImmuno~\cite{DeepImmuno2021}             & DeepImmuno          & prob.                   & both & 9--10 & aware & academic             & \checkmark \\
deepHLApan~\cite{deepHLApan2019}             & deepHLApan (Imm)    & prob.                   & both & 8--14 & aware & academic             & \checkmark \\
PRIME-2.1~\cite{PRIME2023}                   & PRIME               & combined score          & both & 8--14 & aware & academic             & \checkmark \\
PredIG~\cite{PredIG2025}                     & PredIG              & prob.                   & both & 8--14 & aware & academic             & \checkmark \\
ImmuneApp~\cite{ImmuneApp2024}               & Immuno.\ score      & imm.\ score             & both & 8--14 & aware & academic             & \checkmark \\
BigMHC\_IM~\cite{BigMHC2023}                 & BigMHC\_IM          & prob.                   & both & 8--15 & aware & acad.\ (n.c.)        & \checkmark \\
CNNeo$^{\star}$~\cite{CNNeoPP2026}           & CNNeo               & prob.                   & both & 8--14 & aware & MIT                  & \checkmark \\
NeoTImmuML$^{\star}$~\cite{NeoTImmuML2025}   & NeoTImmuML          & prob.                   & both & 8--13 & aware & academic             & \checkmark \\
ICERFIRE-1.0~\cite{ICERFIRE2024}             & icerfire\_score     & $100-$\%rank (DAI)      & MT   & 8--14 & aware & DTU\pendingDTU       & \checkmark \\
NetTepi-1.0~\cite{NetTepi2014}               & MT\_NetTepi         & combined T-cell prop.   & MT   & 8--11 & aware & DTU\pendingDTU       & \checkmark$^{\ddagger}$ \\
MUNIS~\todo{cite MUNIS}                      & MUNIS               & prob.                   & MT   & 8--15 & aware & CC-BY-4.0            & $\circ$ \\
andy90~\todo{cite andy90}                    & amplitude           & prob.                   & MT   & 8--11 & aware & MIT                  & $\circ$ \\
Seq2Neo~\cite{Seq2Neo2022}                   & Seq2Neo\_imm.       & prob.\ (0--1)           & MT   & 8--11 & aware & AFL-3.0\pendingDTU   & $\circ$ \\
ImmugenX~\todo{cite ImmugenX}               & ---                 & prob.                   & MT   & 8--11 & aware & acad.\ SW            & $\circ$ \\
DeepNeo~\todo{cite DeepNeo}                  & ---                 & binding+imm.            & ---  & ---   & aware & ---                  & $\times$ \\
Inference (8-class)~\todo{internal tool}     & class\_0..7         & multiclass              & ---  & ---   & ---   & internal             & $\times$ \\
MHLAPre~\todo{cite MHLAPre}                  & ---                 & prob.                   & ---  & ---   & aware & ---                  & $\times$ \\
\midrule
\multicolumn{8}{l}{\textit{(4) Immunogenicity --- statistical / sequence}}\\
IEDB-Calis~\todo{cite IEDB Calis}           & IEDB\_Calis         & sequence score          & both & 8--11 & \emph{partial} & NPOSL-3.0  & \checkmark \\
Repitope~\todo{cite Repitope}               & Immuno.Score        & sequence score          & both & 8--11 & \emph{agn.} & MIT           & \checkmark \\
\midrule
\multicolumn{8}{l}{\textit{(5) Immunogenicity --- structural / foreignness / pipeline}}\\
NeoaPred~\cite{NeoaPred2024}                 & MT\_NeoaPred        & foreignness (struct.)   & MT   & 9     & aware & Apache-2.0           & $\circ$ \\
T-SCAPE$^{\star}$~\cite{TSCAPE2025}          & MT\_TSCAPE          & T-cell prop.\ (struct.) & MT   & $\le$20 & aware & CC-BY-NC-ND          & \checkmark \\
pTuneos~\cite{pTuneos2019}                   & pTuneos             & pipeline composite      & both & 8--14 & aware & academic             & \checkmark \\
IMPROVE~\cite{IMPROVE2024}                   & IMPROVE             & pipeline composite      & both & 8--12 & aware & academic             & \checkmark \\
ImmunoStruct~\cite{ImmunoStruct2025}         & ---                 & structural              & ---  & ---   & aware & academic             & $\times$ \\
\bottomrule
\end{tabular}%
}
\end{table}

We now state the methodological provenance of the panel, which governs how much
weight each integrated tool's numbers can bear.

\paragraph{Version provenance.} The integrated tools fall into a hierarchy of
decreasing source fidelity, and we annotate this hierarchy explicitly rather than
presenting all numbers as equally authoritative. \emph{(i) Official weights or
official images.} The majority of integrated predictors are run with their
published weights or official containers (for example, MHCflurry, PRIME with
MixMHCpred, ImmuneApp, deepHLApan, PredIG, BigMHC, DeepImmuno, and the structural
tool NeoaPred), and where a vendor reference implementation exists we verified
reproduction against it. \emph{(ii) Honestly degraded configurations.} Two tools
are run in a configuration weaker than their published best, and the limitation
is carried through every table: pTuneos is run through its Pre\&RecNeo sub-model,
the component applicable to our input format; and IMPROVE is run with its
expression feature degraded, since transcript abundance is unavailable for the
benchmark candidates, while the model is otherwise used as published. \emph{(iii)
Self-retrained replicas.} Official pre-trained weights for NeoTImmuML and for the
CNN-based CNNeo are not publicly available; we therefore evaluate models
retrained by us on the published architectures, denoted NeoTImmuML$^{\star}$ and
CNNeo$^{\star}$, and we do not claim that their scores reflect the performance of
the original published models. \emph{(iv) Proxy.} HLAthena is a presentation
predictor used as a presentation \emph{proxy} (Section above), not an
immunogenicity tool. These annotations bound the conclusions we draw about each
tool to the configuration actually evaluated.

\paragraph{HLA-agnostic caveat.} TransHLA and Repitope do not take the HLA allele
as input; under our harmonized protocol they therefore assign the \emph{same}
value to a given peptide across all of a patient's alleles, and IEDB-Calis is
HLA-aware only in part (a position-based score with allele-specific masks for a
subset of alleles). For the fully HLA-agnostic tools, a weak per-patient signal
is the \emph{expected} consequence of their design rather than evidence of poor
ability, and we flag this caveat wherever their results appear; their numbers are
not used as evidence about HLA-aware discrimination.

\paragraph{Restricted-license caveat.} Two distinct kinds of license restriction
apply to subsets of the panel, and we keep them separate because they impose
different obligations. \emph{(a) DTU pending-consent.} netMHCpan (BA, Aff, and
EL), NetMHCstabpan, ICERFIRE, NetTepi, and Seq2Neo (which depends on the DTU tool
netCTLpan) are distributed under DTU academic licenses whose terms forbid
redistributing, without written consent, benchmark numbers computed with their
software. Every such entry is marked \pendingDTU\ in Table~\ref{tab:roster} and
throughout the paper; their numbers are withheld from the main claims and will be
released only once written consent is obtained at the submission stage.
\emph{(b) No-derivatives restriction.} T-SCAPE is distributed under the
CC-BY-NC-ND 4.0 license, whose no-derivatives clause forbids distributing modified
versions of the work. This is a different constraint from DTU pending-consent and
is therefore \emph{not} marked \pendingDTU; it nevertheless requires that the
tool's use be declared and that no modified redistribution occur. None of this
paper's headline conclusions depends on any of the tools in either category.

\subsection{Seven-step harmonization pipeline}
\label{sec:harmonization}

A central methodological contribution is a single pipeline that reduces every
predictor to a comparable peptide-level score, so that observed differences
reflect the tools rather than incidental differences in input handling. The
pipeline proceeds in seven steps:

\begin{enumerate}
  \item \textbf{Unified input.} All candidates are cast to a common record
        format (mutant peptide, source peptide context, and the patient HLA set),
        so that every tool is presented with identical information.
  \item \textbf{Sub-peptide sliding window.} Each peptide is exhaustively
        enumerated into overlapping sub-peptides of lengths $8$--$14$ with a
        stride of $1$, covering the binding-core windows that the predictors
        are designed to score.
  \item \textbf{Sub-model degradation.} For tools that support only a subset of
        lengths or features, the input is restricted accordingly (e.g.,
        DeepImmuno is scored on $9$--$10$-mers only), and the restriction is
        recorded in the coverage table (Table~\ref{tab:coverage}).
  \item \textbf{Re-mapping to the master backbone.} Every sub-peptide~$\times$~HLA
        score is mapped back to its originating master backbone peptide, so that
        scores from different tools are indexed by the same peptide identifier.
  \item \textbf{Unified aggregation.} The per-window, per-allele scores for a
        peptide are aggregated into a single peptide-level value. The primary
        rule is the \emph{best-binder} (maximum) aggregation; \texttt{mean} and
        \texttt{top3mean} are computed as robustness comparators
        (Section~\ref{sec:fairness}).
  \item \textbf{Unified ground truth.} All tools are scored against the same
        ELISpot SFC target for the peptide.
  \item \textbf{Unified evaluation.} A fixed set of thresholds
        (SFC $>0$, $>10$, $>$ median) and a fixed metric set
        (Section~\ref{sec:metrics}) are applied identically to every tool.
\end{enumerate}

The choice of best-binder (maximum-over-windows-and-alleles) aggregation as the
primary rule is grounded in standard neoantigen pipeline practice, in which a
peptide is summarized by its strongest predicted binder across alleles and
lengths~\cite{NetMHCpan2020,PRIME2023}; the same best-binder convention is
codified in established neoantigen pipelines, which summarize a candidate by the
strongest-binding mutant peptide across sub-peptides and alleles (the
``Best MT IC50'' criterion)~\cite{pVACtools2020,pVACseq2016}.
Biologically, this operationalizes the premise that strong presentation by
\emph{any one} of a patient's HLA alleles is sufficient to enable a T-cell
response. We emphasize that, while the maximum rule itself is supported by the
literature, the decision to additionally report \texttt{mean} and
\texttt{top3mean} as sensitivity analyses, to lock the primary comparison to a
single aggregation--threshold combination, and to span the three thresholds
above are our own engineering choices rather than conventions prescribed by any
single source.

\subsection{Per-patient evaluation protocol}
\label{sec:perpatient}

Pooling all peptides across patients and computing a single global Spearman
correlation can mask a tool's ability to rank candidates \emph{within} an
individual patient. We therefore adopt a per-patient protocol that mirrors the
cohort-level evaluation paradigm of prior neoantigen
benchmarks~\cite{TESLA2020,Muller2023}: for each patient $i$ we compute the
Spearman correlation $\rho_i$ between the tool score and ELISpot SFC over that
patient's $n_i$ peptides, and then aggregate the per-patient coefficients.

The primary aggregate is a fixed-effect Fisher-$z$ weighted mean. Each
coefficient is transformed by $z_i = \operatorname{arctanh}(\rho_i)$, with the
Spearman-specific sampling variance
$\operatorname{Var}(z_i) \approx (1+\rho_i^2/2)/(n_i-3)$ and inverse-variance
weight $w_i = (n_i-3)/(1+\rho_i^2/2)$; the pooled estimate is
$\bar{\rho} = \tanh\!\big(\sum_i w_i z_i / \sum_i w_i\big)$, with a $95\%$
confidence interval $\tanh\!\big(\bar{z} \pm 1.96/\sqrt{\sum_i w_i}\big)$. The
median $\tilde{\rho} = \operatorname{median}_i(\rho_i)$ is reported alongside it
as a robust co-primary estimate. The cohort contains only $K=9$ patients with
small per-patient sample sizes ($n_i = 5$--$16$), which is too few to fit a
random-effects model reliably; we therefore report the Fisher-$z$ weighted mean
and its interval as a \emph{descriptive} summary of central tendency rather than
as an estimate of a single common correlation, since the large between-patient
dispersion documented in Section~\ref{sec:res-perpatient} indicates that no common
$\rho$ exists. A random-effects treatment would only widen the intervals, so our
descriptive intervals are, if anything, optimistic. We treat the simple
unweighted mean, the sample-size-weighted (Hunter--Schmidt) mean, the geometric
mean, and the power mean ($p=2$) purely as additional sensitivity analyses; the
last two are computed on the transformed quantity $(1+\rho_i)/2$ and are not used
to construct confidence intervals. We stress that the small $K$ and small $n_i$
imply wide confidence intervals on every per-patient aggregate, and we report
them as such rather than over-interpreting point estimates.

\subsection{Metrics and fairness safeguards}
\label{sec:metrics}

\paragraph{Metrics.} We report three complementary metrics. The area under the
ROC curve (AUC-ROC) measures binary discrimination at each threshold. The area
under the precision--recall curve (AUPRC) is reported as the primary
discrimination metric under class imbalance, where ROC-AUC can be optimistic and
precision--recall analysis is more informative~\cite{Saito2015}. The Spearman
rank correlation between predicted score and ELISpot SFC measures the ability to
recover the continuous magnitude of response; we use a rank correlation rather
than a parametric one because the existing tools emit uncalibrated scores on
heterogeneous scales, so only their ordering relative to magnitude is comparable
across tools. Every reported correlation is accompanied by its $p$-value and,
where applicable, a bootstrap confidence interval. We note that a future tool that
emits a \emph{calibrated} magnitude estimate should additionally be evaluated with
agreement and concordance metrics---such as the concordance index and
Bland--Altman limits of agreement~\cite{Harrell1996,BlandAltman1986}---which the
rank correlation used here does not capture.

\paragraph{Fairness safeguards.}\label{sec:fairness} Five safeguards govern the comparisons.
\emph{(i) Single-criterion lock-down.} The headline ranking is fixed to one
aggregation--threshold combination (best-binder, SFC $>0$); per-tool
``best-of-nine'' selection over aggregation~$\times$~threshold combinations is
explicitly forbidden as a ranking basis, since it would award different tools
different operating points and inflate apparent performance
(\emph{selection-on-max}). \emph{(ii) Transparent coverage.} Tools differ in the
peptides and windows they can score; all coverage gaps and the resulting
effective sample sizes are reported in Table~\ref{tab:coverage} so that
differences in denominators are visible. \emph{(iii) Leakage disclosure.}
Several tools were trained on IEDB-derived data that may overlap with the
benchmark peptides; any such overlap would bias their scores optimistically, and
we disclose this explicitly rather than presenting the tools as fully
independent. \emph{(iv) Proxy separation.} The presentation proxy (HLAthena) is
reported separately and never ranked against immunogenicity predictors.
\emph{(v) Reproducibility check.} The harmonized re-run reproduces the
first-batch scoring to within a maximum absolute AUC difference of $0.004$,
confirming that the pipeline conventions are stable across runs.

\paragraph{Coverage and missingness.} Table~\ref{tab:coverage} summarizes, for
each tool, the number of peptides with a valid aggregated score and the source
of any missing values. Two tools have reduced denominators that we carry through
all downstream analyses: PRIME loses a single peptide (all $210$ of its rows are
unscored because its alleles are unsupported by the underlying predictor),
leaving $n=100$ peptides with $n_{\mathrm{pos}}=89$; deepHLApan loses three
peptides, leaving $n=98$ with $n_{\mathrm{neg}}=10$ rather than $11$. These
differences are small but are reported so that cross-tool comparisons account for
them.

\begin{table}[t]
\centering
\caption{Per-tool coverage and missingness on DS2 ($101$ peptides expanded to
$33{,}922$ sub-peptide~$\times$~HLA rows). ``Valid peptides'' counts peptides
with a non-missing aggregated score. NeoTImmuML$^\star$ denotes the
self-retrained model (official weights unavailable).}
\label{tab:coverage}
\begin{tabular}{lccl}
\toprule
Tool & Valid peptides & NaN rows & Note \\
\midrule
DeepImmuno        & $101$ & $22{,}889/33{,}922$ & $9$--$10$-mers only \\
PredIG            & $101$ & $0$                 & full $8$--$14$-mer coverage \\
IMPROVE           & $101$ & $7{,}457$           & $8$--$12$-mers \\
NeoTImmuML$^\star$ & $101$ & $3{,}508$          & $8$--$13$-mers \\
pTuneos           & $101$ & $0$                 & full coverage \\
PRIME             & $100$ & $210$               & $1$ peptide unsupported (rare allele); $n_{\mathrm{pos}}=89$ \\
ImmuneApp         & $101$ & $0$                 & full coverage \\
deepHLApan        & $98$  & $2{,}065$           & missing across $10$ peptides; $n_{\mathrm{neg}}=10$ \\
\bottomrule
\end{tabular}
\end{table}
